% GENERATED FILE. Edit canonical lessons, then run npm run generate:books.
% canonical-source-hash: fnv1a64:706d4fc611d981e3
% canonical-lessons: ZH-C113-qie, ZH-C113-zhu-boil, ZH-C113-chao, ZH-C113-kao, ZH-C113-jian

\chapter{Cooking}
\label{ch:zh-cooking}
\hlchaptermodality{113}

\begin{chapteropening}
\textbf{By the end of this chapter:} \emph{I can say cut, boil, stir-fry, roast and fry.}

Complete the last lesson of chapter 113: I can say cut, boil, stir-fry, roast and fry.
\end{chapteropening}

\section[qiē]{qiē — to cut; to slice}

\label{lesson:ZH-C113-qie}



\begin{quote}
Before the new one: say the Mandarin for to pull, then the Mandarin for stairs.
\end{quote}

\subsection*{You'll want to know: qiē}
\emph{qiē} — \textquotedblleft{}to cut; to slice\textquotedblright{}.

\begin{quote}
\emph{wǒ qiē cài} — I'll chop the vegetables.
\end{quote}

Its characters come later; for now it is a word for the ear.

\begin{cousinweb}
\emph{qiē} is inherited from Middle Chinese \emph{tshet} (Guangyun, in Baxter's transcription), glossed 'to cut' and 'to carve'.
\end{cousinweb}

\begin{grammarlens}[title={What you've built}]
One more word for everyday actions.
\end{grammarlens}

\subsection*{Your turn}
\begin{itemize}
\raggedright
  \item \emph{Say it:} \emph{qiē}
  \item \emph{Say it:} \emph{qiē}, once more
  \item \emph{Say it:} \emph{wǒ qiē cài} — I'll chop the vegetables.
  \item \emph{Recall:} say the Mandarin for to pull, then the Mandarin for stairs, then say \emph{qiē} again
\end{itemize}

\begin{tcolorbox}[breakable,colback=teal!4,colframe=teal!35!black,title={Before you move on}]
What does qiē mean? (\textquotedblleft{}To cut; to slice\textquotedblright{}.) Say it once more.
\end{tcolorbox}
\section[zhǔ]{zhǔ — to boil; to cook in water}

\label{lesson:ZH-C113-zhu-boil}



\begin{quote}
Before the new one: say the Mandarin for stairs, then the Mandarin for to slice.
\end{quote}

\subsection*{You'll want to know: zhǔ}
\emph{zhǔ} — \textquotedblleft{}to boil; to cook in water\textquotedblright{}.

\begin{quote}
\emph{wǒmen zhǔ jiǎozi} — We're boiling dumplings.
\end{quote}

Its characters come later; for now it is a word for the ear.

\begin{cousinweb}
\emph{zhǔ} is inherited from Middle Chinese \emph{tsyoX} (Guangyun, in Baxter's transcription), listed under a variant form of the character and glossed 'to cook'.
\end{cousinweb}

\begin{grammarlens}[title={What you've built}]
One more word for everyday actions.
\end{grammarlens}

\subsection*{Your turn}
\begin{itemize}
\raggedright
  \item \emph{Say it:} \emph{zhǔ}
  \item \emph{Say it:} \emph{zhǔ}, once more
  \item \emph{Say it:} \emph{wǒmen zhǔ jiǎozi} — We're boiling dumplings.
  \item \emph{Recall:} say the Mandarin for stairs, then the Mandarin for to slice, then say \emph{zhǔ} again
\end{itemize}

\begin{tcolorbox}[breakable,colback=teal!4,colframe=teal!35!black,title={Before you move on}]
What does zhǔ mean? (\textquotedblleft{}To boil; to cook in water\textquotedblright{}.) Say it once more.
\end{tcolorbox}
\section[chǎo]{chǎo — to stir-fry}

\label{lesson:ZH-C113-chao}



\begin{quote}
Before the new one: say the Mandarin for to slice, then the Mandarin for to boil.
\end{quote}

\subsection*{You'll want to know: chǎo}
\emph{chǎo} — \textquotedblleft{}to stir-fry\textquotedblright{}.

\begin{quote}
\emph{nǐ huì chǎo cài ma} — Can you stir-fry vegetables?
\end{quote}

Its characters come later; for now it is a word for the ear.

\begin{cousinweb}
\emph{chǎo} is inherited from Middle Chinese \emph{tsrhaewX} (Guangyun, in Baxter's transcription), glossed 'to roast dry'.
\end{cousinweb}

\begin{grammarlens}[title={What you've built}]
One more word for everyday actions.
\end{grammarlens}

\subsection*{Your turn}
\begin{itemize}
\raggedright
  \item \emph{Say it:} \emph{chǎo}
  \item \emph{Say it:} \emph{chǎo}, once more
  \item \emph{Say it:} \emph{nǐ huì chǎo cài ma} — Can you stir-fry vegetables?
  \item \emph{Recall:} say the Mandarin for to slice, then the Mandarin for to boil, then say \emph{chǎo} again
\end{itemize}

\begin{tcolorbox}[breakable,colback=teal!4,colframe=teal!35!black,title={Before you move on}]
What does chǎo mean? (\textquotedblleft{}To stir-fry\textquotedblright{}.) Say it once more.
\end{tcolorbox}
\section[kǎo]{kǎo — to roast; to bake}

\label{lesson:ZH-C113-kao}



\begin{quote}
Before the new one: say the Mandarin for to boil, then the Mandarin for to stir-fry.
\end{quote}

\subsection*{You'll want to know: kǎo}
\emph{kǎo} — \textquotedblleft{}to roast; to bake\textquotedblright{}.

\begin{quote}
\emph{wǒmen kǎo miànbāo} — We're making toast.
\end{quote}

Its characters come later; for now it is a word for the ear.

\begin{cousinweb}
\emph{kǎo} is the everyday word for roasting and baking. No dictionary we use records an older form, so its history is not settled.
\end{cousinweb}

\begin{grammarlens}[title={What you've built}]
One more word for everyday actions.
\end{grammarlens}

\subsection*{Your turn}
\begin{itemize}
\raggedright
  \item \emph{Say it:} \emph{kǎo}
  \item \emph{Say it:} \emph{kǎo}, once more
  \item \emph{Say it:} \emph{wǒmen kǎo miànbāo} — We're making toast.
  \item \emph{Recall:} say the Mandarin for to boil, then the Mandarin for to stir-fry, then say \emph{kǎo} again
\end{itemize}

\begin{tcolorbox}[breakable,colback=teal!4,colframe=teal!35!black,title={Before you move on}]
What does kǎo mean? (\textquotedblleft{}To roast; to bake\textquotedblright{}.) Say it once more.
\end{tcolorbox}
\section[jiān]{jiān — to fry (in a little oil)}

\label{lesson:ZH-C113-jian}



\begin{quote}
Before the new one: say the Mandarin for to stir-fry, then the Mandarin for to roast.
\end{quote}

\subsection*{You'll want to know: jiān}
\emph{jiān} — \textquotedblleft{}to fry (in a little oil)\textquotedblright{}.

\begin{quote}
\emph{wǒ huì jiān jīdàn} — I can fry an egg.
\end{quote}

Its characters come later; for now it is a word for the ear.

\begin{cousinweb}
\emph{jiān} is inherited from Middle Chinese \emph{tsjen} (Guangyun, in Baxter's transcription), glossed 'to cook until done'. Frying in a little oil is the sense today; the step is not given there.
\end{cousinweb}

\begin{grammarlens}[title={What you've built}]
One more word for everyday actions.
\end{grammarlens}

\subsection*{Your turn}
\begin{itemize}
\raggedright
  \item \emph{Say it:} \emph{jiān}
  \item \emph{Say it:} \emph{jiān}, once more
  \item \emph{Say it:} \emph{wǒ huì jiān jīdàn} — I can fry an egg.
  \item \emph{Recall:} say the Mandarin for to stir-fry, then the Mandarin for to roast, then say \emph{jiān} again
\end{itemize}

\begin{tcolorbox}[breakable,colback=teal!4,colframe=teal!35!black,title={Before you move on}]
What does jiān mean? (\textquotedblleft{}To fry (in a little oil)\textquotedblright{}.) Say it once more.
\end{tcolorbox}
